% !TEX root = ../../main.tex

\section{系统异常处理与用户反馈实现}

HearBridge 系统由移动端、后台管理端、Spring Boot 服务端和 Python 手势识别服务协同组成，功能链路中涉及接口调用、文件上传、实时通信、模型加载和语义修正等多个环节。
因此，系统需要对常见异常进行统一处理，并将底层错误转换为用户能够理解的提示信息。

在移动端场景中，常见异常包括登录状态失效、网络连接失败、相机权限不可用、WebSocket 连接中断、未检测到有效动作、视频文件选择失败和视频识别失败等。
对于这些情况，移动端应通过页面提示、加载状态、错误弹窗或结果区域提示等方式向用户说明当前状态，避免用户停留在无反馈界面。

在后台管理端场景中，常见异常包括管理员登录失效、资源字段缺失、资源编码重复、文件上传失败、样本删除失败、训练任务触发失败和模型发布失败等。
对于新增、修改、删除和发布等关键操作，后台管理端应提供必要的确认提示，并在操作完成后展示成功或失败反馈。
对于模型发布这类会影响识别服务运行状态的操作，系统应特别注意数据库状态与识别服务实际加载状态的一致性。

在服务端实现中，异常处理需要覆盖参数校验、权限校验、数据库访问、对象存储访问、外部服务调用和识别服务调用等情况。
服务端不应直接将底层异常暴露给前端，而应转换为统一响应结构，包含状态码、错误信息和必要的业务说明。
这样可以降低前端处理复杂度，也可以避免用户看到难以理解的系统错误。

在句子视频识别链路中，可能出现视频文件上传失败、视频格式不支持、Python 识别服务不可用、模型未加载、语义修正服务调用失败等问题。
服务端需要将这些错误转换为移动端可理解的提示信息。
若语义修正失败但基础识别成功，系统仍可以返回原始识别结果和中文映射结果，并提示语义修正暂不可用。
这样可以避免单个外部服务异常导致整个识别流程完全不可用。

通过异常处理和用户反馈机制，系统能够在复杂链路下保持较好的可用性和可理解性。
用户侧可以明确知道当前操作是否成功、失败原因是什么以及是否可以重新尝试；
管理员侧可以根据提示判断资源维护、样本处理或模型发布是否完成。
该机制有助于提升系统整体稳定性和使用体验。
